\chapter{Edge Colourings}
%%%%%%%%%%%%%%%%%%%%%%%%%%%%%%%%%%%%%%%%%%%%%%%%%%%%%%%%%%%%%%%%%%%%%%%%%%%%%%%
\section{Edge Colourings}

% Generated by scripts/blueprint_restate.py from TCSlib.GraphTheory.Core.EdgeColourings.
% Each body is the STATEMENT only, taken from the declaration's Lean docstring:
% the one-line summary plus the verbatim book statement.  Proof, proof plan,
% significance commentary and formalisation notes are excluded by construction.

\begin{definition}[Colour $i$ is \emph{represented} at $v$: some edge incident with $v$ has colour $i$]
\lean{SimpleGraph.IsRepresentedAt}
\label{SimpleGraph.IsRepresentedAt}
\leanok
Colour \texttt{i} is \emph{represented} at \texttt{v}: some edge incident with
\texttt{v} has colour \texttt{i}. A \texttt{k}-edge colouring here is a plain function
\texttt{C : G.edgeSet $\to$ Fin k} --- deliberately \textbf{not} proper.

A $k$-\emph{edge colouring} $\mathscr{C}$ of a loopless graph $G$ is an assignment of
$k$ colours, $1, 2, \ldots, k$, to the edges of $G$. The colouring $\mathscr{C}$ is
\emph{proper} if no two adjacent edges have the same colour.

We say that colour $i$ is \emph{represented} at vertex $v$ if some edge incident with
$v$ has colour $i$.
\end{definition}

\begin{definition}[$c(v)$: the number of distinct colours represented at $v$]
\lean{SimpleGraph.numColoursAt}
\label{SimpleGraph.numColoursAt}
\leanok
\uses{SimpleGraph.IsRepresentedAt}
\texttt{c(v)}: the number of distinct colours represented at \texttt{v}.

Given a $k$-edge colouring $\mathscr{C}$ of $G$ we shall denote by $c(v)$ the number of
distinct colours represented at $v$. Clearly, we always have \[c(v) \le d(v) \qquad
(6.3)\] Moreover, $\mathscr{C}$ is a proper $k$-edge colouring if and only if equality
holds in (6.3) for all vertices $v$ of $G$.
\end{definition}

\begin{definition}[An optimal $k$-edge colouring: one that cannot be improved, where an improvement strictly increases $\sum\dots$]
\lean{SimpleGraph.IsOptimalEdgeColouring}
\label{SimpleGraph.IsOptimalEdgeColouring}
\leanok
\uses{SimpleGraph.numColoursAt}
An \textbf{optimal} \texttt{k}-edge colouring: one that cannot be improved, where an
improvement strictly increases \texttt{$\sum$ v, c(v)}. A static extremal condition, not
a procedure.

We shall call a $k$-edge colouring $\mathscr{C}'$ an \emph{improvement} on $\mathscr{C}$
if \[\sum_{v \in V} c'(v) > \sum_{v \in V} c(v)\] where $c'(v)$ is the number of
distinct colours represented at $v$ in the colouring $\mathscr{C}'$. An \emph{optimal}
$k$-edge colouring is one which cannot be improved.
\end{definition}

\begin{definition}[$G(E_i \cup E_j)$, the subgraph on the edges coloured $i$ or $j$]
\lean{SimpleGraph.twoColourSubgraph}
\label{SimpleGraph.twoColourSubgraph}
\leanok
\texttt{G[E\_i $\cup$ E\_j]}, the subgraph on the edges coloured \texttt{i} or \texttt{j}.
\end{definition}

\begin{definition}[$G$ is uniquely $k$-edge-colourable: any two proper $k$-edge colourings agree up to a permutation of the\dots]
\lean{SimpleGraph.IsUniquelyEdgeColourable}
\label{SimpleGraph.IsUniquelyEdgeColourable}
\leanok
\texttt{G} is \textbf{uniquely \texttt{k}-edge-colourable}: any two proper \texttt{k}-edge colourings agree up to a
permutation of the colours.

\textbf{6.2.5} $G$ is called \emph{uniquely $k$-edge-colourable} if any two proper
$k$-edge colourings of $G$ induce the same partition of $E$.
\end{definition}

\begin{definition}[The Petersen graph]
\lean{SimpleGraph.petersenGraph}
\label{SimpleGraph.petersenGraph}
The Petersen graph.
\end{definition}

\begin{theorem}[Maximum degree bounds the edge chromatic number]
\lean{SimpleGraph.maxDegree_le_edgeChromaticNumber}
\label{SimpleGraph.maxDegree_le_edgeChromaticNumber}
Let $G$ be a finite simple graph, and let $L(G)$ denote its line graph — the graph whose
vertices are the edges of $G$, with two such vertices adjacent exactly when the
corresponding edges of $G$ share an endpoint. Then the maximum degree $\Delta(G)$ of $G$
is at most the chromatic number of $L(G)$; that is, writing $\chi'(G) = \chi(L(G))$ for
the edge chromatic number of $G$, one has \[\chi'(G) \geq \Delta(G).\]
\end{theorem}

\begin{theorem}[The number of colours at a vertex is at most its degree]
\lean{SimpleGraph.numColoursAt_le_degree}
\label{SimpleGraph.numColoursAt_le_degree}
\uses{SimpleGraph.numColoursAt}
Let $G$ be a finite loopless graph, and let $\mathscr{C}$ be a $k$-edge colouring of
$G$, that is, an arbitrary assignment of one of the colours $1, 2, \ldots, k$ to each
edge (not necessarily proper). For a vertex $v$, write $c(v)$ for the number of distinct
colours represented at $v$, and $d(v)$ for its degree. Then \[c(v) \le d(v).\]
\end{theorem}

\begin{theorem}[Properness in terms of colours at each vertex]
\lean{SimpleGraph.isProper_iff_numColoursAt_eq_degree}
\label{SimpleGraph.isProper_iff_numColoursAt_eq_degree}
\uses{SimpleGraph.numColoursAt}
Let $G$ be a finite graph and let $C$ be a $k$-edge colouring of $G$, that is, an
assignment of one of the colours $1, 2, \ldots, k$ to each edge. Then $C$ is proper — no
two distinct edges that share a common endpoint receive the same colour — if and only if
$c(v) = d(v)$ for every vertex $v$, where $c(v)$ is the number of distinct colours
represented at $v$ and $d(v)$ is the degree of $v$.
\end{theorem}

\begin{theorem}[A two-colouring of edges with both colours at every vertex of degree ≥ 2]
\lean{SimpleGraph.exists_two_edge_colouring_both_represented}
\label{SimpleGraph.exists_two_edge_colouring_both_represented}
\uses{SimpleGraph.IsRepresentedAt}
Let $G$ be a finite connected simple graph that is not an odd cycle (that is, $G$ is not
a single cycle of odd length using every edge of $G$). Then the edges of $G$ can be
assigned two colours in such a way that at every vertex $v$ of degree at least $2$, both
colours are represented at $v$ — each of the two colours appears on some edge incident
with $v$.
\end{theorem}

\begin{theorem}[Optimal edge colourings force odd-cycle components]
\lean{SimpleGraph.isOddCycle_component_of_isOptimalEdgeColouring}
\label{SimpleGraph.isOddCycle_component_of_isOptimalEdgeColouring}
\uses{SimpleGraph.IsOptimalEdgeColouring, SimpleGraph.IsRepresentedAt, SimpleGraph.twoColourSubgraph}
Let $\mathscr{C}$ be an optimal $k$-edge colouring of a finite graph $G$, and suppose
there is a vertex $u$ and colours $i, j$ for which colour $i$ is not represented at $u$
while at least two edges incident with $u$ carry colour $j$. Then the connected
component of $u$ in the subgraph $G[E_i \cup E_j]$ spanned by the edges coloured $i$ or
$j$ is an odd cycle: there is a closed walk based at $u$ that is a cycle of odd length
whose vertices are precisely those of that component.
\end{theorem}

\begin{theorem}[König's edge-colouring theorem for bipartite graphs]
\lean{SimpleGraph.edgeChromaticNumber_eq_maxDegree_of_isBipartite}
\label{SimpleGraph.edgeChromaticNumber_eq_maxDegree_of_isBipartite}
Let $G$ be a bipartite simple graph on a finite vertex set, and let $\Delta(G)$ denote
its maximum degree. Then the chromatic number of the line graph of $G$ — equivalently,
the least number of colours needed to colour the edges of $G$ so that no two edges
sharing an endpoint receive the same colour — equals $\Delta(G)$.
\end{theorem}

\begin{theorem}[Vizing's theorem on the chromatic index]
\lean{SimpleGraph.vizing_chromatic_index}
\label{SimpleGraph.vizing_chromatic_index}
Let $G$ be a simple graph on a finite vertex set, with maximum degree $\Delta(G)$, and
let $\chi'(G)$ denote its chromatic index — that is, the chromatic number of the line
graph of $G$, equivalently the least number of colours needed to colour the edges of $G$
so that any two edges sharing an endpoint receive distinct colours. Then either
$\chi'(G) = \Delta(G)$ or $\chi'(G) = \Delta(G) + 1$.
\end{theorem}

\begin{theorem}[Vizing's theorem: chromatic index at most one more than maximum degree]
\lean{SimpleGraph.vizing_chromatic_index_le}
\label{SimpleGraph.vizing_chromatic_index_le}
Let $G$ be a finite simple graph with maximum degree $\Delta(G)$, and let $L(G)$ denote
its line graph. Then the chromatic number of $L(G)$ — equivalently, the chromatic index
(edge chromatic number) $\chi'(G)$ of $G$ — satisfies
\[
\chi(L(G)) \le \Delta(G) + 1.
\]
\end{theorem}

\begin{theorem}[Rebalancing two disjoint matchings]
\lean{SimpleGraph.exists_rebalanced_matchings}
\label{SimpleGraph.exists_rebalanced_matchings}
Let $G$ be a finite simple graph, and let $M$ and $N$ be two disjoint matchings of $G$
(each a set of edges, no two of which share a vertex) with $\abs{N} < \abs{M}$. Then
there exist disjoint matchings $M'$ and $N'$ of $G$ such that $\abs{M'} = \abs{M} - 1$,
$\abs{N'} = \abs{N} + 1$, and $M' \cup N' = M \cup N$.
\end{theorem}

\begin{theorem}[Balanced matching decomposition of a bipartite graph]
\lean{SimpleGraph.exists_balanced_matching_decomposition}
\label{SimpleGraph.exists_balanced_matching_decomposition}
Let $G$ be a finite bipartite graph with edge set $E$ of size $\varepsilon = |E|$ and
maximum degree $\Delta$, and let $p$ be a positive integer with $\Delta \le p$. Then
there exist $p$ matchings $M_1, M_2, \ldots, M_p$ of $G$ that are pairwise edge-disjoint
and satisfy
\[ E = M_1 \cup M_2 \cup \cdots \cup M_p, \]
with each of their sizes controlled by
\[ \left\lfloor \varepsilon/p \right\rfloor \le |M_i| \le \left\lceil \varepsilon/p \right\rceil \qquad (1 \le i \le p); \]
in particular, any two of the matchings differ in size by at most one.
\end{theorem}

\begin{theorem}[Regular bipartite supergraph of a bipartite graph]
\lean{SimpleGraph.exists_regular_bipartite_supergraph}
\label{SimpleGraph.exists_regular_bipartite_supergraph}
Let $G$ be a finite bipartite simple graph, and let $\Delta$ denote its maximum degree.
Then there is a finite bipartite simple graph $H$ that is $\Delta$-regular (every vertex
of $H$ has degree exactly $\Delta$) together with an embedding of $G$ into $H$ — an
injection $V(G) \to V(H)$ under which two vertices of $G$ are adjacent if and only if
their images are adjacent in $H$. In other words, every finite bipartite graph is
(isomorphic to) an induced subgraph of a bipartite graph regular of degree equal to
$\Delta$.
\end{theorem}

\begin{theorem}[Every colour represented at each vertex in a bipartite graph]
\lean{SimpleGraph.exists_edge_colouring_all_represented_of_isBipartite}
\label{SimpleGraph.exists_edge_colouring_all_represented_of_isBipartite}
\uses{SimpleGraph.IsRepresentedAt}
Let $G$ be a finite bipartite simple graph with minimum degree $\delta > 0$. Then there
is an edge colouring $C$ of $G$ using $\delta$ colours — that is, a function assigning
to each edge one of the colours $1, 2, \ldots, \delta$, not required to be proper — such
that every colour is represented at every vertex: for each vertex $v$ and each colour
$i$, some edge incident with $v$ is assigned colour $i$.
\end{theorem}

\begin{theorem}[Edge-chromatic number of odd complete graphs]
\lean{SimpleGraph.edgeChromaticNumber_completeGraph_odd}
\label{SimpleGraph.edgeChromaticNumber_completeGraph_odd}
For every integer $n \ge 1$, the complete graph $K_{2n-1}$ on $2n-1$ vertices has
edge-chromatic number equal to $2n-1$; that is, its edges can be properly coloured with
$2n-1$ colours but no fewer. Equivalently, the chromatic number of the line graph of
$K_{2n-1}$ is $2n-1$.
\end{theorem}

\begin{theorem}[Edge chromatic number of the even complete graph]
\lean{SimpleGraph.edgeChromaticNumber_completeGraph_even}
\label{SimpleGraph.edgeChromaticNumber_completeGraph_even}
For every integer $n \ge 1$, the complete graph $K_{2n}$ on $2n$ vertices has edge
chromatic number $\chi'(K_{2n}) = 2n - 1$; equivalently, the chromatic number of the
line graph of $K_{2n}$ equals $2n - 1$.
\end{theorem}

\begin{theorem}[Edge chromatic number of an odd-order regular graph]
\lean{SimpleGraph.edgeChromaticNumber_eq_maxDegree_add_one_of_regular_odd_card}
\label{SimpleGraph.edgeChromaticNumber_eq_maxDegree_add_one_of_regular_odd_card}
Let $G$ be a finite simple graph whose number of vertices is odd, and suppose $G$ is
$k$-regular for some integer $k \ge 1$ (every vertex has exactly $k$ neighbours). Then
the edge chromatic number of $G$ — equivalently, the chromatic number of its line graph
— equals $\Delta(G) + 1$, where $\Delta(G)$ is the maximum degree of $G$. Writing
$\chi'(G)$ for the edge chromatic number, \[\chi'(G) = \Delta(G) + 1.\]
\end{theorem}

\begin{theorem}[Odd graphs with many edges are class two]
\lean{SimpleGraph.edgeChromaticNumber_eq_maxDegree_add_one_of_card_edges_gt}
\label{SimpleGraph.edgeChromaticNumber_eq_maxDegree_add_one_of_card_edges_gt}
Let $G$ be a finite simple graph whose vertex set has $2n+1$ elements for some $n \in
\bbn$, and write $\Delta$ for its maximum degree and $\varepsilon$ for its number of
edges. If $\varepsilon > n\Delta$, then the edge chromatic number of $G$ equals $\Delta
+ 1$; that is, the least number of colours needed to colour the edges of $G$ so that no
two edges sharing an endpoint receive the same colour is exactly $\Delta + 1$.
\end{theorem}

\begin{theorem}[Chromatic index of an even regular graph with one edge subdivided]
\lean{SimpleGraph.edgeChromaticNumber_subdivide_eq_maxDegree_add_one}
\label{SimpleGraph.edgeChromaticNumber_subdivide_eq_maxDegree_add_one}
\uses{SimpleGraph.subdivide}
Let $G$ be a $k$-regular simple graph, with $k \ge 2$, on a finite vertex set whose
cardinality is even, and let $uv$ be an edge of $G$. Form the graph $G'$ by subdividing
the edge $uv$: delete it and join its ends $u$ and $v$ by a path of length two through a
single new vertex. Then the chromatic number of the line graph of $G'$ — equivalently,
the edge chromatic number $\chi'(G')$ — equals $\Delta(G') + 1$, where $\Delta(G')$ is
the maximum degree of $G'$.
\end{theorem}

\begin{theorem}[Chromatic index of a lightly edge-deleted odd-order regular graph]
\lean{SimpleGraph.edgeChromaticNumber_deleteEdges_eq_maxDegree_add_one}
\label{SimpleGraph.edgeChromaticNumber_deleteEdges_eq_maxDegree_add_one}
Let $G$ be a finite simple $k$-regular graph with $k \ge 1$ whose vertex set has odd
cardinality, and let $F$ be a set of edges of $G$ satisfying $2\abs{F} < k$ (that is,
fewer than $k/2$ edges). Write $G - F$ for the graph obtained from $G$ by deleting the
edges in $F$, and let $\Delta$ denote its maximum degree. Then the chromatic index of $G
- F$—equivalently, the chromatic number of its line graph—equals $\Delta + 1$. (Beineke
and Wilson.)
\end{theorem}

\begin{theorem}[Uniquely 3-edge-colourable cubic graphs are Hamiltonian]
\lean{SimpleGraph.isHamiltonian_of_uniquely_three_edge_colourable}
\label{SimpleGraph.isHamiltonian_of_uniquely_three_edge_colourable}
\uses{SimpleGraph.IsUniquelyEdgeColourable}
Let $G$ be a finite simple graph that is $3$-regular and admits a proper
$3$-edge-colouring. If $G$ is uniquely $3$-edge-colourable — meaning any two proper
$3$-edge-colourings of $G$ agree up to a permutation of the three colours — then $G$
contains a Hamiltonian cycle, that is, a closed walk that visits every vertex exactly
once.
\end{theorem}

\begin{theorem}[Chromatic index of the prism over a graph]
\lean{SimpleGraph.edgeChromaticNumber_boxProd_completeGraph_two}
\label{SimpleGraph.edgeChromaticNumber_boxProd_completeGraph_two}
Let $G$ be a finite simple graph, let $K_2$ denote the complete graph on two vertices,
and form the Cartesian product $G \square K_2$ — the prism over $G$, whose vertices are
the pairs $(v,i)$ with $v$ a vertex of $G$ and $i \in \{0,1\}$, and in which $(v,i)$ is
adjacent to $(v',i')$ exactly when either $v = v'$ and $i \neq i'$, or $i = i'$ and
$v,v'$ are adjacent in $G$. Then the edge chromatic number of $G \square K_2$ — that is,
the chromatic number of its line graph — equals its maximum degree $\Delta(G \square
K_2)$.
\end{theorem}

\begin{theorem}[Chromatic index of a Cartesian product of graphs]
\lean{SimpleGraph.edgeChromaticNumber_boxProd}
\label{SimpleGraph.edgeChromaticNumber_boxProd}
Let $G$ and $H$ be finite simple graphs, and let $G \square H$ denote their Cartesian
product, the graph on $V(G) \times V(H)$ in which $(u,x)$ and $(v,y)$ are adjacent
exactly when either $u = v$ and $x \sim y$ in $H$, or $x = y$ and $u \sim v$ in $G$. For
a graph $K$, write $\Delta(K)$ for its maximum degree and $\chi'(K)$ for its chromatic
index, the least number of colors needed to color the edges of $K$ so that any two edges
meeting at a common vertex receive different colors (equivalently, the chromatic number
of the line graph of $K$). If $H$ has at least one edge, so that $\Delta(H) > 0$, and
$H$ satisfies $\chi'(H) = \Delta(H)$, then $\chi'(G \square H) = \Delta(G \square H)$.
\end{theorem}
